\chapter{Судалгааны үр дүн}
\label{ch:results}

\section{Дуу танилтын харьцуулалт}
\label{sec:res-stt}

13 бичлэгийн нэг нь Gemini-д хязгаарын алдаа (HTTP 429) өгсөн. Иймээс хоёр систем хоёулаа хариу өгсөн 12 бичлэгээр харьцуулав. Эдгээрт нийт 176 үг байна. Үр дүнг \tref{tab:stt_results}-т харуулав.

\begin{table}[H]
\centering
\caption{Ижил 12 бичлэгийн үр дүн}
\label{tab:stt_results}
\begin{tabular}{|l|r|r|}
\hline
\textbf{Хэмжүүр} & \textbf{Whisper} & \textbf{Gemini 3.5 Transcribe} \\
\hline
Үгийн алдаа & 67 / 176 & 71 / 176 \\
WER & 38.07\% & 40.34\% \\
\hline
CER & 17.93\% & 24.85\% \\
\hline
Нэг бичлэгийн хугацаа (медиан) & 3.98 сек & 6.96 сек \\
\hline
\end{tabular}
\end{table}

Whisper 67, Gemini 71 үг алдсан. Ялгаа 4 үг тул их биш. Тэмдэгтийн алдаа Gemini-д илүү их гарсан. Мөн нэг бичлэгт хэлсэн «2026» оныг Gemini «2020» гэж таньсан, Whisper зөв таньсан. Сургуулийн мэдээлэл оноос их хамаардаг тул энэ алдаа чухал. Gemini урьдчилсан шалгуурт тэнцээгүй. Gemini-ийн хугацаанд файл илгээх, хариу авах хугацаа орсон.

Бүх бичлэг нэг хүнийх, тоо нь цөөн. Тиймээс аль арга нь ерөнхийдөө илүү гэж хэлэх боломжгүй.

\section{Сонголт}
\label{sec:res-choice}

Одоогоор чатад Gemini 3.5 Transcribe-ийг ашиглаж байна. Учир нь яриаг текст болгох тусгай үйлчилгээ, компьютерт загвар ажиллуулах шаардлагагүй, үнэгүй квоттой. Гэхдээ интернэт, API түлхүүр хэрэгтэй, аудио гадагш явдаг. Whisper нь компьютер дээрээ ажилладаг ч Монгол хэлэнд сайжруулахын тулд сургах хэрэгтэй \cite{mengke2025khalkha}.

Эцсийн сонголтыг шинэ бичлэгээр дахин харьцуулж, удирдагч багштайгаа ярилцаж хийнэ.

% Дараа нэмэх: \section{Хайлт ба хариултын үр дүн} — W07–W08-ын хэмжилтийн дараа
